%% ma: L12-B19-0002
%% lop: 12
%% chuong: 6
%% bai: 19
%% dang: 12.19.2
%% loai: TLN
%% muc: VD
%% dap_an: "0,08"
%% nguon: "0. ĐỀ THAM KHẢO TN THPT 2025.pdf (D00), Phần III Câu 6"
%% kien_thuc: "Bài 19 (công thức xác suất toàn phần, công thức Bayes), Bài 18 (xác suất có điều kiện)"
%% trang_thai: chinh-thuc
\begin{ex}
Có hai chiếc hộp, hộp I có $6$ quả bóng màu đỏ và $4$ quả bóng màu vàng, hộp II có $7$ quả bóng màu đỏ và $3$ quả bóng màu vàng, các quả bóng có cùng kích thước và khối lượng. Lấy ngẫu nhiên một quả bóng từ hộp I bỏ vào hộp II. Sau đó, lấy ra ngẫu nhiên một quả bóng từ hộp II. Tính xác suất để quả bóng được lấy ra từ hộp II là quả bóng được chuyển từ hộp I sang, biết rằng quả bóng đó có màu đỏ (làm tròn kết quả đến hàng phần trăm).
\shortans{0,08}
\loigiai{Hộp II sau khi nhận thêm một quả có $11$ quả. Gọi $A$: ``quả lấy ra từ hộp II màu đỏ'', $B$: ``quả lấy ra là quả chuyển từ hộp I sang''.
$P(A)=\frac{6}{10}\cdot\frac{8}{11}+\frac{4}{10}\cdot\frac{7}{11}=\frac{76}{110}$, $P(A\cap B)=\frac{6}{10}\cdot\frac{1}{11}=\frac{6}{110}$.
Vậy $P(B\mid A)=\dfrac{P(A\cap B)}{P(A)}=\dfrac{6}{76}\approx0{,}08$.}
\end{ex}
